\scanpage{70}
% archive foliation: 69
$i=0$ \unsure{seulmt}) $C^i\xrightarrow{d^i}C^{i+1}$ ait un Ker \deleted{\unsure{qui soit} tel que $Z^i\longrightarrow C^i$ est un $\Sigma$-mono (Ici, le $\Sigma$ dans $B$ est de la plus \unsure{mauvaise}, il n'y a \unsure{pas à prouver}]).} Ceci dit, on se borne à de tels complexes, qui correspondent donc à des foncteurs \deleted{\unsure{de $B$}} $k$-addit. \inserted{($\Sigma$-ex. à g)} $\unsure{\mathcal P|P}_0\longrightarrow B$ resp. $Q\unsure{\longrightarrow|{-}}B$.

\deleted{On se propose de \inserted{se} donner une structure multiplicative sur ceux-ci, ce qui nous \uline{oblige} à donner de plus
\[
\deleted{F(\Psi^0)\otimes_B F(C^i)\longrightarrow F(\Psi^0*C^i)\quad\unsure{\text{et de plus}}}
\]
un hom de $B$, \unsure{à foncteurs|de bifoncteurs}
\[
F(X)\otimes_B F(Y)\longrightarrow F(X*Y)
\]
sur $\unsure{\mathcal P|P}_0\times \unsure{\mathcal P|P}_0$ \inserted{$\unsure{\mathcal P|P}$} resp $Q\times Q$. Or}

\unsure{Si notre classe $B$ est stable à} \unsure{un noyau}, alors on se ramène aux foncteurs $\Sigma$-ex. à g de $\unsure{\mathcal P|P}$ dans $B$, en \deleted{prenant} \unsure{interprétant} ceux-ci comme des foncteurs $k$-additifs quelc. sur $K$ i.e. des $C^\bullet\in\operatorname{Ob}B^\bullet$. On peut l'expliciter ainsi: $\forall L\in \unsure{\mathcal P|P}$ \inserted{($\in Q$, $\in \unsure{\mathcal P|P}_0$)} \deleted{\unsure{on sait}} s'insère dans deux $\Sigma$-suites de $\unsure{\mathcal P|P}$ ($Q,\unsure{\mathcal P|P}_0$)
\[
\unsure{0}{-}L{-}I{-}L'{-}\unsure{0}
\]
\[
0{-}L'\longrightarrow J{-}L''\longrightarrow0
\]
avec $I,J\in K$, et par suite on

